\chapter{Conclusion}

This thesis investigated graph-aware prompting for entity typing using local WiKC taxonomy structure. The completed study compared standard, chain-of-thought, and contrastive prompting with and without combined graph evidence, added parent-, sibling-, and child-only conditions, and conducted a matched Llama seed-0 child-removal experiment. The supplied result archives contain 37,300 completed outputs from 37,353 planned entity–condition–run combinations. Incomplete results and attempts are retained and reported.


RQ1 receives a consistent answer within the experimental scope: graph-aware prompting improves F1, precision, recall, MRR, and depth-weighted nDCG across Gemini and all three Llama seeds. Paired F1 gains range from 0.0451 to 0.1348. The treatment supplies relevant oracle structure together with interpretation instructions, so it does not isolate evidence alone.


RQ2 shows that parent-only context outperforms sibling-only context on F1, contrary to H1. Including child context does not produce a general F1 gain in the matched seed-0 experiment. RQ3 shows gains generally shrinking with depth, contrary to H2. RQ4 shows higher Gemini F1 against each Llama seed in all nine matched condition comparisons, while the three Llama seeds expose variation in the relative strength of relation-only prompts.


The evaluation distinguishes several dimensions of performance. Recovering a maximum-depth gold label early is different from recovering an entire ancestor chain. Hierarchical closeness is different from explicitly listing a closed hierarchy. Exact benchmark agreement is different from complete semantic validity. These distinctions explain why the thesis presents several metrics, paired effects, distributions, and inspectable cases instead of one headline accuracy value.


The development process helped define that interpretation. Exploratory prompts were refined, an insufficiently useful taxonomy-extraction idea was abandoned, and generation controls were fixed for the final external-demonstration runs. Only the supplied new response archives contribute experimental scores; earlier outputs are excluded.


The empty-child control suggests a template contribution: point estimates range from about +0.005 to +0.024, with Gemini near +0.02 compared with its +0.0625 full-cohort standard gain. These different-cohort effects cannot be subtracted to isolate evidence, and the seed-123 interval includes zero. Gemini thinking-token increases also co-vary with graph gains, so structural evidence and additional computation are not isolated.


The conclusion is specific: oracle-selected local taxonomy evidence and its interpretation instructions form a useful inference-time prompt treatment for this ranked multi-label typing task. The study provides a reproducible basis for testing retrieved evidence, taxonomy quality, controlled noise, and independent semantic judgments.

